% 380_D4_Spezifitaet_De_ch.tex
% Johann Pascher, 29. Sept. 2026
% Warum D4: Spezifität des Trägers gegen angepasste Vergleichsgitter
% (offene Brücke R95, Dok. 190)

\providecommand{\BM}{\textbf{[B]}}
\providecommand{\KM}{\textbf{[K]}}
\providecommand{\SM}{\textbf{[S]}}
\providecommand{\QM}{\textbf{[Q]}}
\providecommand{\SetzM}{\textbf{[SETZUNG]}}

\chapter*{Warum \texorpdfstring{$D_4$}{D4}}
\addcontentsline{toc}{chapter}{Warum D4}

\begin{abstract}
Die FFGFT setzt das Gitter $D_4$ als Träger der kompakten Geometrie
$T^4/\mathbb{Z}_3$ (Dok.~314) \SetzM. Dok.~190 führt unter R95 als offene
Brücke, dass ein adversarieller Spezifitätstest gegen angepasste
Vergleichsträger fehlt. Dieses Dokument prüft den Träger. Drei vorab
festgelegte Kriterien werden auf $\mathbb{Z}^4$, $A_4$, $A_2\oplus A_2$, $D_4$ und auf
Zufallsgitter angewandt: die im Korpus geführte fixpunktfreie
$\mathbb{Z}_3$-Wirkung mit neun Fixpunkten (Dok.~330), die Kusszahl und die
Gitterenergie als Epstein-Zeta-Funktion bei gleichem Volumen der
Grundmasche. Nur das dritte Kriterium ist von der FFGFT unabhängig. Die
$\mathbb{Z}_3$-Bedingung schließt $\mathbb{Z}^4$ und $A_4$ aus; unter den
verbleibenden Trägern hat $D_4$ die kleinste Gitterenergie \KM, und die
Kusszahl 24 ist mit Dok.~340 konsistent. $D_4$ ist damit als Träger besser
motiviert, bleibt aber eine Setzung; R95 ist nicht geschlossen, die Ebene
des Operators $\hat F$ bleibt offen \SM. Prüfskript: 13/13.
\end{abstract}

% ============================================================
\section*{Statusmarkierungen}
% ============================================================
\begin{center}
\begin{tabular}{@{}lp{10cm}@{}}
\textbf{[B]} & Algebraisch bewiesen. \\[3pt]
\textbf{[K]} & Numerisch bestätigt (Prüfskript). \\[3pt]
\textbf{[S]} & Strukturell plausibel, offen. \\[3pt]
\textbf{[Q]} & Aus der Literatur übernommen, zitiert. \\[3pt]
\textbf{[SETZUNG]} & Axiomatische Setzung. \\
\end{tabular}
\end{center}

% ============================================================
\section{Die Frage}
% ============================================================

Eine Setzung ist motiviert, aber nicht beliebig, wenn der gesetzte Träger
unter vergleichbaren Alternativen durch Größen hervorsticht, die nicht an die
Ergebnisse angepasst wurden. Genau das verlangt die GAE-Logik, die im
HLV-Audit angewandt wurde: Dort übertraf ein gewöhnliches kubisches Gitter den
untersuchten Träger im Spektralmaß, und die Spezifitätsbehauptung fiel
(Dok.~330). Für $D_4$ hält R95 fest, dass $5\nmid|\mathrm{Aut}(D_4)|=1152$
eine gitterspezifische Tatsache ist, aber kein Spezifitätsbeweis.

Der Test muss daher drei Bedingungen erfüllen. Die Vergleichsträger haben
dieselbe Dimension und dasselbe Volumen der Grundmasche. Die Kriterien sind
vor der Rechnung festgelegt. Und ein Kriterium, das Spezifität belegen soll,
darf keine Zahl sein, die erst aus der FFGFT stammt. Kriterien, die eine
Korpusbedingung abbilden, prüfen dagegen nur die Verträglichkeit; sie werden
im Folgenden ausdrücklich so gekennzeichnet.

% ============================================================
\section{Vergleichsträger und Kriterien}
% ============================================================

Verglichen werden das kubische Gitter $\mathbb{Z}^4$ (der Gewinner im
HLV-Audit), das Wurzelgitter $A_4$, die Summe $A_2\oplus A_2$ und $D_4$,
jeweils auf das Volumen~1 der Grundmasche skaliert, dazu Zufallsgitter.

\paragraph{K1: $\mathbb{Z}_3$-Verträglichkeit (Korpusbedingung).}
Der Korpus führt $T^4/\mathbb{Z}_3$ mit neun Fixpunkten und einer
$\mathbb{Z}_3$-Wirkung mit den Eigenwerten $\{\omega,\omega^2,\omega,\omega^2\}$;
die neun Fixpunkte folgen aus $|\det(1-g)|=|(1-\omega)(1-\omega^2)|^2=9$
(Dok.~330, ebenso Dok.~361, 364, 365). Gefordert ist also ein
Gitterautomorphismus $g$ der Ordnung~3 ohne Eigenwert~1. K1 bildet diese
Bedingung ab und prüft Verträglichkeit, nicht Spezifität. Die Suche ist
vollständig, weil jeder Automorphismus kürzeste Vektoren auf kürzeste
Vektoren abbildet und alle vier Basen aus kürzesten Vektoren bestehen.

Welche $\mathbb{Z}_3$ gemeint ist, zeigt die vollständige Aufzählung der 1152
Automorphismen von $D_4$: 80 haben die Ordnung~3. Davon besitzen 64 einen
zweidimensionalen Fixraum und erzeugen Fixtori statt isolierter Fixpunkte,
darunter der Dynkin-Automorphismus der Trialität. Die übrigen 16 sind
fixpunktfrei mit genau neun Fixpunkten auf $T^4$ \KM. Das sind die
16~Elemente, die Dok.~330 nennt; die Trialität, die das Orbifold trägt, ist
also diese fixpunktfreie Klasse und nicht der Dynkin-Automorphismus.

\paragraph{K2: Kusszahl (Konsistenz).}
Die Zahl der kürzesten Gittervektoren. Sie ist eine Eigenschaft des Gitters
allein, aber die FFGFT verwendet sie selbst ($\mathrm{Kissing}(D_4)=24$ in
Dok.~340). K2 prüft daher Konsistenz, nicht Spezifität.

\paragraph{K3: Gitterenergie.}
Die Epstein-Zeta-Funktion
\begin{equation}
Z_\Lambda(\nu)=\sum_{v\in\Lambda\setminus\{0\}}|v|^{-\nu}
\end{equation}
bei Volumen~1 der Grundmasche, für $\nu=5,6,8$. Sie ist zugleich die
spektrale Zeta-Funktion des flachen Torus zum dualen Gitter, also eine
Operatorgröße des Laplace-Operators, und enthält keinen freien Parameter.
Berechnet wird sie mit EpsteinLib (Buchheit et al.) \QM; die Bibliothek
reproduziert die geschlossenen Formeln für $\mathbb{Z}^4$ und $D_4$ auf zwölf
Stellen \KM.

% ============================================================
\section{Ergebnisse}
% ============================================================

\paragraph{K1.}
$\mathbb{Z}^4$ und $A_4$ besitzen keinen Automorphismus der Ordnung~3 ohne
Eigenwert~1; auf ihnen ist $T^4/\mathbb{Z}_3$ mit neun Fixpunkten nicht
möglich \KM. $A_2\oplus A_2$ und $D_4$ besitzen einen solchen, jeweils mit
$|\det(1-g)|=9$ \KM. Der Gewinner des HLV-Audits scheidet also schon an der
Symmetriebedingung aus.

\paragraph{K2.}
Die Kusszahlen sind $D_4$: 24, $A_4$: 20, $A_2\oplus A_2$: 12,
$\mathbb{Z}^4$: 8 \KM. Unter den beiden $\mathbb{Z}_3$-verträglichen Trägern
hat nur $D_4$ die 24, die in Dok.~340 als $\mathrm{Kissing}(D_4)$ eingeht.
Das ist Konsistenz: Der Korpus rechnet mit dem Gitter, das er setzt.

\paragraph{K3.}
\begin{center}
\begin{tabular}{lcccc}
\toprule
$\nu$ & $D_4$ & $A_4$ & $A_2\oplus A_2$ & $\mathbb{Z}^4$ \\
\midrule
5 & \textbf{22,860} & 23,123 & 23,703 & 24,531 \\
6 & \textbf{12,583} & 12,919 & 13,699 & 14,830 \\
8 & \textbf{6,830} & 7,282 & 8,457 & 10,246 \\
\bottomrule
\end{tabular}
\end{center}
$D_4$ hat bei allen drei Exponenten die kleinste Energie \KM. Von 300
allgemeinen Zufallsgittern unterschreitet keines $D_4$ (bestes $13{,}15$ bei
$\nu=6$), ebenso keines von 300 Zufallsgittern aus der
$\mathbb{Z}_3$-verträglichen Familie (bestes $12{,}72$) \KM. Alle 200 kleinen
Störungen um $D_4$ erhöhen die Energie; $D_4$ ist ein lokales Minimum \KM.

Bewiesen ist, dass $D_4$ ein lokales Minimum der Epstein-Zeta-Funktion unter
allen vierdimensionalen Gittern ist (Sarnak und Strömbergsson 2006) \QM. Ob
es auch das globale Minimum ist, ist in vier Dimensionen eine offene
Vermutung. Die 600 Zufallsgitter stützen sie, beweisen sie aber nicht.

% ============================================================
\section{Einordnung}
% ============================================================

\paragraph{Was gezeigt ist.}
Unter den vierdimensionalen Trägern, die die im Korpus geführte
$\mathbb{Z}_3$-Wirkung mit neun Fixpunkten zulassen, hat $D_4$ die kleinste
Gitterenergie \KM. Das ist die einzige der drei Größen, die nicht aus der
FFGFT stammt; die Kusszahl bestätigt nur die Konsistenz. $D_4$ ist damit als
Träger besser motiviert als bisher, bleibt aber eine motivierte Setzung
\SetzM. Der Befund ersetzt keinen Spezifitätsbeweis.

\paragraph{Was offen bleibt.}
Die GAE-Logik fragt auch nach dem Operator: ob $\hat F$ mit
$\xi=\lambda_{\min}(\hat F_{D_4})$ (Dok.~327) auf den Vergleichsträgern
etwas Unterscheidbares liefert. Das ist hier nicht geprüft \SM. Ebenso ist
die Zufallssuche ein numerischer Beleg, kein Beweis der globalen Minimalität
in der $\mathbb{Z}_3$-Familie; sie deckt die Exponenten $\nu=5,6,8$ ab.

\paragraph{Folgen für die Buchführung.}
R95 ist nicht geschlossen. Der geforderte Operatortest steht weiter aus, und
neue Resultate bekommen keinen Registereintrag.

% ============================================================
\section{Zusammenfassung}
% ============================================================

\begin{center}
\begin{tabular}{>{\raggedright\arraybackslash}p{3.2cm}cccc}
\toprule
Kriterium & $D_4$ & $A_2\oplus A_2$ & $A_4$ & $\mathbb{Z}^4$ \\
\midrule
K1: $\mathbb{Z}_3$ mit 9 Fixpunkten (Korpus) & ja & ja & nein & nein \\
K2: Kusszahl (Konsistenz) & \textbf{24} & 12 & 20 & 8 \\
K3: $Z(6)$ bei Volumen 1 & \textbf{12,58} & 13,70 & 12,92 & 14,83 \\
\bottomrule
\end{tabular}
\end{center}

Unter den Trägern, die die Symmetriebedingung erfüllen, hat $D_4$ die
kleinste Gitterenergie, eine Größe, die nicht aus der FFGFT stammt; die
Kusszahl ist mit dem Korpus konsistent. Die Setzung des Trägers ist damit
besser gestützt, bleibt aber eine Setzung.

\paragraph{Prüfskript.} \texttt{2/python/Dok380\_Skripte/pruef\_380\_d4\_spezifitaet.py}
(13/13): Kontrolle der Bibliothek an geschlossenen Formeln, K1 bis K3,
Klassifikation der Ordnung-3-Automorphismen von $D_4$, Zufalls- und
Störungstests.

\paragraph{Quellen.}
Buchheit,~A.~A.; Busse,~J.~K.; Gutendorf,~R.: \textit{Computation and
Properties of the Epstein Zeta Function with Applications to Quantum
Systems.} arXiv:2412.16317 (2024); Software EpsteinLib,
\url{https://github.com/epsteinlib/epsteinlib}.
Sarnak,~P.; Strömbergsson,~A.: \textit{Minima of Epstein's zeta function and
heights of flat tori.} Invent.\ Math.\ \textbf{165}, 115--151 (2006).
Conway,~J.~H.; Sloane,~N.~J.~A.: \textit{Sphere Packings, Lattices and
Groups.} 3.~Aufl., Springer (1999).
